% !TeX root = ../../Book5.tex
% * 状态：初稿已完成；正文、证明、算例与配套习题解答已写齐，尚未正式审核。
% 纲要标题：## 附录 G　仿射 Lie 代数初步

\chapter{仿射 Lie 代数初步}
\label{b5:app:G}
\BookChapterNumber{b5:app:G}

\begin{chapterabstract}
	有限型 Cartan 矩阵的正定性失去一个方向后，可以得到仿射型矩阵。相应的无限维结构在 Laurent 多项式环上有具体模型：环 Lie 代数加入一个中心方向，再加入记录次数的导子。本附录证明这些构造及其根空间分解，以仿射 \(\mathfrak{sl}_2\) 完整建立最高权与可积性的判据，并通过形式特征标比较不同级的权空间。Weyl–Kac 公式给出证明概要，秩一、秩二的低次计算则直接展开。
\end{chapterabstract}

\begin{learninggoals}
	\begin{itemize}
		\item 从广义 Cartan 矩阵构造实现，区分正定有限型与余秩一半正定仿射型，并计算二阶例子。
		\item 验证环 Lie 代数的中心余圈及次数导子，求中心、导出代数和实根、虚根的重数。
		\item 在仿射 \(\mathfrak{sl}_2\) 中核验生成关系，构造 Verma 模与简单最高权模，并证明可积性的整数条件。
		\item 计算级一最高权模的低次空间，解释每条简单根串有限与整个模无限维为何能够同时成立。
		\item 在明确的形式完成中使用 Weyl–Kac 公式，辨认分母中的虚根重数，并核验小秩特征标系数。
	\end{itemize}
\end{learninggoals}

把 \(\mathfrak{sl}_2(\mathbb C)\) 的每个元素附上一个整数次数，就得到 \(x\otimes t^m\) 这样的向量，括号先按 \([x\otimes t^m,y\otimes t^n]=[x,y]\otimes t^{m+n}\) 计算。这个无限维代数已有无穷多个模态，但定义级所需的中心元尚未出现。加入中心项以后，公式变成
\[
	[h\otimes t^m,h\otimes t^{-m}]=2m c.
\]
若中心元 \(c\) 作用为非零标量，正、负次数的 Cartan 算子便不能同时作用在一个有限维空间上：它们的交换子的迹应为零，而右边的迹不会为零。这是无限维表示在本附录中自然出现的一个原因。

\ProseParagraphBreak
Cartan 矩阵也能先于具体模型定义代数。Kac 与 Moody\footnote{卡茨（\OriginalName{Виктор Гершевич Кац}，Victor G. Kac，1943-），俄裔美国数学家，研究无限维 Lie 代数及其表示；穆迪（Robert Vaughan Moody，1941-），加拿大数学家，研究 Lie 理论与非周期序。两人独立发展了后来称为 Kac--Moody 代数的理论。}的构造保留简单根生成元及 Serre 关系，允许矩阵不再正定。这里先说明有限型与仿射型的差别，再把仿射秩一关系放进环模型核验，建立关系代数的三角分解，计算环模型的括号与根空间，并证明秩一最高权模的可积性。

\ProseParagraphBreak
所需基础是\chapref{b5:ch:14}的 \(\mathfrak{sl}_2\) 表示、\chapref{b5:ch:15}的 PBW 定理、\chapref{b5:ch:16}的根空间分解，以及\chapref{b5:ch:17}的 Cartan 矩阵与生成关系。\chapref{b5:ch:18}的最高权构造在无限维环境下重新核对；\chapref{b5:ch:19}的有限 Weyl 特征标公式提供比较。底域固定为 \(\mathbb C\)，实际表示计算集中在未扭仿射模型。

\ProseParagraphBreak
一般仿射展示与环实现的同构识别，以及特征标理论，采用 Kac\cite[Section 2, pp. 300--301]{Kac1980HighestWeight}的结果。该文先用 Cartan 型代数的结构描述环中心扩张，再以 Verma 模反变型及简单子商的限制证明特征标公式。本附录说明这些结果在证明概要中的作用，并展开无限乘积的系数解释、仿射反射与低次特征标的计算。

% !TeX root = ../../Book5.tex
% 纲要标题：### G.1 广义 Cartan 矩阵的有限型与仿射型
% 纲要位置：工作记录/计划纲要.md（2026-10-20 新增骨架）。
% 写作状态：2026-10-20 新增骨架；正文、证明、算例与习题均待撰写。
% 写作范围：以可对称化及仿射型为范围，展开环代数和仿射 sl2 小模型，Weyl–Kac 公式标明证明范围；不展开 VOA 与 Lie 超代数。
% #### G.1.1 广义 Cartan 矩阵
% #### G.1.2 有限型与仿射型
% #### G.1.3 生成关系与可对称化条件

\section{广义 Cartan 矩阵的有限型与仿射型}
\label{b5:sec:G01}

待撰写。

\subsection{广义 Cartan 矩阵}
\label{b5:subsec:G01:01}

待撰写。

\subsection{有限型与仿射型}
\label{b5:subsec:G01:02}

待撰写。

\subsection{生成关系与可对称化条件}
\label{b5:subsec:G01:03}

待撰写。

% 新增结构与纲要同步；以下条目均为待撰写计划。
% * 状态：待撰写；当前仅有结构占位。

% !TeX root = ../../Book5.tex
% 纲要标题：### G.2 环 Lie 代数、中心扩张与导子
% 纲要位置：工作记录/计划纲要.md（2026-10-20 新增骨架）。
% 写作状态：2026-10-20 新增骨架；正文、证明、算例与习题均待撰写。
% 写作范围：以可对称化及仿射型为范围，展开环代数和仿射 sl2 小模型，Weyl–Kac 公式标明证明范围；不展开 VOA 与 Lie 超代数。
% #### G.2.1 环 Lie 代数
% #### G.2.2 中心扩张与二次型
% #### G.2.3 次数导子与根分解

\section{环 Lie 代数、中心扩张与导子}
\label{b5:sec:G02}

待撰写。

\subsection{环 Lie 代数}
\label{b5:subsec:G02:01}

待撰写。

\subsection{中心扩张与二次型}
\label{b5:subsec:G02:02}

待撰写。

\subsection{次数导子与根分解}
\label{b5:subsec:G02:03}

待撰写。

% 新增结构与纲要同步；以下条目均为待撰写计划。
% * 状态：待撰写；当前仅有结构占位。

% !TeX root = ../../Book5.tex
% * 状态：初稿已完成；正文、证明、算例与配套习题解答已写齐，尚未正式审核。
% 纲要标题：### G.3 仿射秩一代数、级与可积最高权
% #### G.3.1 仿射秩一代数的生成关系
% #### G.3.2 级与最高权
% #### G.3.3 可积条件与秩一算例
% 写作范围：在未扭仿射 sl2 的环实现中证明生成、最高权构造及可积判据。

\section{仿射秩一代数、级与可积最高权}
\label{b5:sec:G03}

有限维 \(\mathfrak{sl}_2\) 表示中的最高权必须是非负整数，这个整数决定了下降链的长度。仿射代数含有两套相互作用的 \(\mathfrak{sl}_2\) 三元组，因此最高权要同时满足两个整数条件。不过，每一条简单根方向的链都有限，并不意味着整个模有限维：两种下降操作可以交替进行，经过无限多个不同的权。下面把这两点放到同一个环代数模型中计算。

\subsection{仿射秩一代数的生成关系}
\label{b5:subsec:G03:01}

取 \(\mathfrak{sl}_2(\mathbb C)\) 的标准基 \(e,f,h\)，满足 \([h,e]=2e\)、\([h,f]=-2f\)、\([e,f]=h\)。在上一节用迹型构造的
\[
	\widehat{\mathfrak{sl}}_2
	=\bigl(\mathfrak{sl}_2(\mathbb C)\otimes\mathbb C[t,t^{-1}]\bigr)
	\oplus\mathbb Cc\oplus\mathbb C\mathrm d
\]
中，定义
\begin{equation}\label{b5:eq:G-affine-chevalley-elements}
	\begin{aligned}
		e_1&=e\otimes1,& f_1&=f\otimes1,& h_1&=h,\\
		e_0&=f\otimes t,& f_0&=e\otimes t^{-1},& h_0&=c-h.
	\end{aligned}
\end{equation}
这里将 \(h\otimes1\) 简写为 \(h\)。下标 \(0,1\) 标记两条简单根，因而 \(e_0\) 并不表示零次模态。第二套三元组把正、负根向量对调，并分别配上 \(t\) 和 \(t^{-1}\)。中心项恰好使它们的括号成为 \(c-h\)。为简洁起见，本节把上一节在 \(\mathfrak{sl}_2\) 情形中的正、负子代数 \(\widehat{\mathfrak n}_\pm\) 简记为 \(\mathfrak n_\pm\)。

\begin{proposition}\label{b5:prop:G-affine-sl2-generators}
	设 \(\widehat{\mathfrak{sl}}_2\) 为未扭复仿射 Lie 代数，其环实现的中心扩张使用 \((x,y)=\operatorname{tr}(xy)\)，取\eqref{b5:eq:G-affine-chevalley-elements}中的元素。令
	\[
		A=(a_{ij})_{i,j=0,1}=\begin{pmatrix}2&-2\\-2&2\end{pmatrix}.
	\]
	则 \(e_0,e_1,f_0,f_1,h_0,h_1,\mathrm d\) 生成 \(\widehat{\mathfrak{sl}}_2\)，并满足
	\[
		\begin{gathered}
			[h_i,h_j]=0,\qquad [h_i,e_j]=a_{ij}e_j,\qquad
			[h_i,f_j]=-a_{ij}f_j,\qquad [e_i,f_j]=\delta_{ij}h_i,\\
			(\operatorname{ad}e_i)^3e_j=(\operatorname{ad}f_i)^3f_j=0
			\quad(i\ne j),\\
			[\mathrm d,h_i]=0,\qquad [\mathrm d,e_i]=\delta_{i0}e_i,
			\qquad [\mathrm d,f_i]=-\delta_{i0}f_i.
		\end{gathered}
	\]
	此外，正、负子代数
	\[
		\begin{aligned}
			\mathfrak n_+&=\mathbb C(e\otimes1)\oplus
			\bigl(\mathfrak{sl}_2(\mathbb C)\otimes t\mathbb C[t]\bigr),\\
			\mathfrak n_-&=\mathbb C(f\otimes1)\oplus
			\bigl(\mathfrak{sl}_2(\mathbb C)\otimes t^{-1}\mathbb C[t^{-1}]\bigr)
		\end{aligned}
	\]
	分别由 \(e_0,e_1\) 和 \(f_0,f_1\) 生成。
\end{proposition}
\begin{proof}
	由 \(\operatorname{tr}(fe)=1\)，有
	\[
		[e_0,f_0]=[f\otimes t,e\otimes t^{-1}]=-h+c=h_0.
	\]
	不同下标的 \(e_i,f_j\) 括号为零；\(h_0=c-h\)、\(h_1=h\) 给出矩阵 \(A\) 的四个系数。次数导子的关系直接由模态次数得到。三次关系可以沿括号计算，例如
	\[
		[e_1,e_0]=h\otimes t,\qquad
		[e_1,[e_1,e_0]]=-2e\otimes t,\qquad
		[e_1,[e_1,[e_1,e_0]]]=0.
	\]
	交换 \(0,1\) 后，前两次括号分别为 \(-h\otimes t\)、\(-2f\otimes t^2\)，第三次仍为零。将 \(e,f,t\) 换成 \(f,e,t^{-1}\)，得到两条负生成元的三次关系。

	为说明生成性，先由 \([e_1,e_0]=h\otimes t\) 得到次数一的 Cartan 方向。利用
	\[
		[h\otimes t,e\otimes t^r]=2e\otimes t^{r+1},\qquad
		[h\otimes t,f\otimes t^r]=-2f\otimes t^{r+1},
	\]
	从 \(e\otimes1\) 和 \(f\otimes t\) 逐次得到全部 \(e\otimes t^r\)（\(r\ge0\)）和 \(f\otimes t^r\)（\(r\ge1\)）；再用 \([e\otimes1,f\otimes t^r]=h\otimes t^r\) 得到全部正次 Cartan 模态。这证明 \(e_0,e_1\) 生成 \(\mathfrak n_+\)。负方向从 \([f_1,f_0]=-h\otimes t^{-1}\) 出发，同样逐次得到 \(\mathfrak n_-\)。最后 \(h_1=h\)、\(h_0+h_1=c\)，再加入 \(\mathrm d\)，便得到整个代数。
\end{proof}

这个计算把广义 Cartan 矩阵的关系落实到了具体模态。验证关系并证明这些元素生成，只给出从关系定义的代数到环实现的满同态；要断言两者同构，还须排除关系代数中的额外核。本节的表示构造直接在上述环实现中进行。

\subsection{级与最高权}
\label{b5:subsec:G03:02}

记 \(\widehat{\mathfrak h}=\mathbb Ch\oplus\mathbb Cc\oplus\mathbb C\mathrm d\)。上一节的根分解给出简单根 \(\alpha_1\) 与 \(\alpha_0=\delta-\alpha_1\)，其中
\[
	\alpha_1(h)=2,\quad \alpha_1(c)=\alpha_1(\mathrm d)=0,
	\qquad \delta(h)=\delta(c)=0,\quad\delta(\mathrm d)=1.
\]
两者作为 \(\widehat{\mathfrak h}^*\) 的元素线性无关。定义 \(\Lambda_0,\Lambda_1\in\widehat{\mathfrak h}^*\) 为
\[
	\Lambda_i(h_j)=\delta_{ij},\qquad \Lambda_i(\mathrm d)=0
	\quad(i,j\in\{0,1\}).
\]
它们是当前实现中的\term[fangshe-jiben-quan]{仿射基本权}[affine fundamental weight]。因为 \(h_0,h_1,\mathrm d\) 是 \(\widehat{\mathfrak h}\) 的基，每个权唯一写成
\[
	\lambda=m_0\Lambda_0+m_1\Lambda_1+a\delta,
	\qquad m_i=\lambda(h_i),\quad a=\lambda(\mathrm d).
\]
这里 \(m_i\) 暂时可以是任意复数，尚未施加整性条件。
\symidx{\(\Lambda_0,\Lambda_1\)：仿射 \(\mathfrak{sl}_2\) 的基本权}

\begin{definition}\label{b5:def:G-level-highest}
	设 \(V\) 为复 Lie 代数 \(\widehat{\mathfrak{sl}}_2\) 的模。如果中心元 \(c\) 在 \(V\) 上按标量 \(\ell\in\mathbb C\) 作用，就称 \(V\) 的\term[ji-level]{级}[level]为 \(\ell\)。给定 \(\lambda\in\widehat{\mathfrak h}^*\)，若存在非零向量 \(v\) 满足
	\[
		Hv=\lambda(H)v\quad(H\in\widehat{\mathfrak h}),\qquad
		e_0v=e_1v=0,\qquad V=U(\widehat{\mathfrak{sl}}_2)v,
	\]
	就称 \(V\) 为最高权 \(\lambda\) 的模，\(v\) 为最高权生成元。
\end{definition}
\symidx{\(\ell=\lambda(c)\)：仿射最高权模的级}
由于 \(e_0,e_1\) 生成 \(\mathfrak n_+\)，定义中的两个零化条件等价于 \(\mathfrak n_+v=0\)。又因为 \(c\) 中心化整个代数，它在由 \(v\) 生成的模上恒按 \(\lambda(c)\) 作用。因此最高权模自动有确定的级，而且
\begin{equation}\label{b5:eq:G-level-sum}
	\ell=\lambda(c)=\lambda(h_0+h_1)=m_0+m_1.
\end{equation}
级与 \(\mathrm d\) 的特征值不同：前者在整个模上固定，后者随着正、负次模态的作用而改变。

\ProseParagraphBreak

有限维半单情形的 Verma 模构造只用到三角分解与 PBW。这里仍有向量空间直和
\[
	\widehat{\mathfrak{sl}}_2
	=\mathfrak n_-\oplus\widehat{\mathfrak h}\oplus\mathfrak n_+,
	\qquad \mathfrak b=\widehat{\mathfrak h}\oplus\mathfrak n_+.
\]
令 \(\mathbb C_\lambda\) 为一维 \(\mathfrak b\) 模，其中 \(\widehat{\mathfrak h}\) 按 \(\lambda\) 作用，\(\mathfrak n_+\) 作用为零；由于 \([\mathfrak b,\mathfrak b]\subseteq\mathfrak n_+\)，这是良定义的作用。相应的仿射 Verma 模是
\[
	M(\lambda)=U(\widehat{\mathfrak{sl}}_2)
	\otimes_{U(\mathfrak b)}\mathbb C_\lambda,
	\qquad v_\lambda=1\otimes1.
\]
\symidx{\(M(\lambda),L(\lambda)\)：仿射 Verma 模及其唯一简单商}

\begin{proposition}\label{b5:prop:G-affine-verma}
	设 \(\lambda\in\widehat{\mathfrak h}^*\)，相对于复 Lie 代数 \(\widehat{\mathfrak{sl}}_2\) 的上述三角分解构造 \(M(\lambda)\)。则 \(u\mapsto uv_\lambda\) 给出向量空间同构 \(U(\mathfrak n_-)\cong M(\lambda)\)。模 \(M(\lambda)\) 按 \(\widehat{\mathfrak h}\) 的权分解，各权空间有限维；其权都属于
	\[
		\lambda-Q_+,\qquad Q_+=\mathbb Z_{\ge0}\alpha_0+\mathbb Z_{\ge0}\alpha_1,
	\]
	且最高权空间为 \(\mathbb Cv_\lambda\)。每个最高权为 \(\lambda\) 的模都是 \(M(\lambda)\) 的商。模 \(M(\lambda)\) 有唯一极大真子模，从而有唯一简单商 \(L(\lambda)\)；不同最高权的简单模互不同构。
\end{proposition}
\begin{proof}
	\cref{b5:thm:pbw}允许 Lie 代数具有任意维数，其证明中的每个字与每次括号展开仍然只含有限项。将 \(\mathfrak n_-\) 的基排在 \(\mathfrak b\) 的基之前，便得到右 \(U(\mathfrak b)\) 模同构
	\[
		U(\mathfrak n_-)\otimes_{\mathbb C}U(\mathfrak b)
		\longrightarrow U(\widehat{\mathfrak{sl}}_2).
	\]
	张量 \(\mathbb C_\lambda\) 后得到第一个结论，特别地 \(v_\lambda\ne0\)。由于 \(\mathfrak n_-\) 由 \(f_0,f_1\) 生成，\(M(\lambda)\) 由 \(f_{i_1}\cdots f_{i_r}v_\lambda\) 张成；其中每个 \(f_i\) 将权降低 \(\alpha_i\)。权 \(\lambda-r_0\alpha_0-r_1\alpha_1\) 只需考虑含 \(r_0\) 个 \(f_0\) 和 \(r_1\) 个 \(f_1\) 的有限多个字，故该权空间有限维。\(\alpha_0,\alpha_1\) 线性无关，因此只有空字能有权 \(\lambda\)。最高权向量 \(v\) 满足定义诱导模所需的关系，故 \(u\otimes1\mapsto uv\) 给出到它所生成模的唯一满同态。

	任意子模都按权分解。具体地，一个向量只含有限多个不同权分量，可选 \(H\in\widehat{\mathfrak h}\) 区分这些权，再用 \(H\) 的插值多项式逐个提取权分量；子模对这些多项式保持不变。每个真子模的 \(\lambda\) 权空间必为零，否则它含有生成元 \(v_\lambda\)。于是所有真子模之和仍没有 \(\lambda\) 权分量，是唯一极大真子模。取商得到 \(L(\lambda)\)。若 \(L(\lambda)\cong L(\mu)\)，两个最高权生成元在同一简单模中都生成全模，所以 \(\mu\in\lambda-Q_+\) 且 \(\lambda\in\mu-Q_+\)；简单根线性无关使两式只能在 \(\lambda=\mu\) 时同时成立。
\end{proof}

\subsection{可积条件与秩一算例}
\label{b5:subsec:G03:03}

\begin{definition}\label{b5:def:G-integrable-module}
	设 \(V\) 为复 Lie 代数 \(\widehat{\mathfrak{sl}}_2\) 的模。若 \(V\) 是 \(\widehat{\mathfrak h}\) 权空间的直和，且对每个 \(v\in V\) 及 \(i\in\{0,1\}\)，都存在正整数 \(N\) 使
	\[
		e_i^Nv=f_i^Nv=0,
	\]
	就称 \(V\) 为\term[keji-mo]{可积模}[integrable module]。这里允许 \(N\) 随 \(v\) 和 \(i\) 改变，这个条件称为相应算子的局部幂零性。
\end{definition}
由两个 \(\mathfrak{sl}_2\) 子代数得到的有限根串，正是可积条件逐向量控制的对象：每个向量沿简单根方向的升降链有限，所需的幂零指数可以随向量改变，整个模仍可有无穷多个权。

\begin{theorem}\label{b5:thm:G-affine-integrable}
	设 \(\lambda\in\widehat{\mathfrak h}^*\)，\(L(\lambda)\) 为复仿射 Lie 代数 \(\widehat{\mathfrak{sl}}_2\) 的简单最高权模，\(v_\lambda\) 为其最高权生成元。则 \(L(\lambda)\) 可积，当且仅当
	\[
		m_0=\lambda(h_0)\in\mathbb Z_{\ge0},\qquad
		m_1=\lambda(h_1)\in\mathbb Z_{\ge0}.
	\]
	在此情形中，\(f_i^{m_i+1}v_\lambda=0\)，而 \(v_\lambda,f_iv_\lambda,\ldots,f_i^{m_i}v_\lambda\) 都非零。特别地，可积简单最高权模的级为非负整数 \(m_0+m_1\)。
\end{theorem}
\begin{proof}
	先固定 \(i\)。三元组 \(e_i,f_i,h_i\) 的括号关系给出与\eqref{b5:eq:sl2-verma-actions}相同的恒等式
	\begin{equation}\label{b5:eq:G-simple-root-string}
		e_i f_i^r v_\lambda=r(m_i-r+1)f_i^{r-1}v_\lambda.
	\end{equation}
	若 \(L(\lambda)\) 可积，取 \(N\ge0\) 使 \(f_i^Nv_\lambda\ne0\)、\(f_i^{N+1}v_\lambda=0\)。在最后一个等式上作用 \(e_i\)，得到 \((N+1)(m_i-N)f_i^Nv_\lambda=0\)，故 \(m_i=N\) 为非负整数。这也给出所述根串长度。

	反过来，设两个 \(m_i\) 都为非负整数。在 \(L(\lambda)\) 中令 \(w_i=f_i^{m_i+1}v_\lambda\)。由\eqref{b5:eq:G-simple-root-string}，\(e_iw_i=0\)；对于 \(j\ne i\)，\([e_j,f_i]=0\) 给出 \(e_jw_i=0\)。若 \(w_i\ne0\)，它是权 \(\lambda-(m_i+1)\alpha_i\) 的最高权向量。由三角分解，它生成的子模所有权都在 \(\lambda-(m_i+1)\alpha_i-Q_+\) 中，不可能含有权 \(\lambda\)，与 \(L(\lambda)\) 的简单性矛盾。因此 \(w_i=0\)。

	还须把生成元上的幂零关系推广到整个模。对固定的 \(i\)，导子 \(\operatorname{ad}f_i\) 在 Lie 代数生成元上局部幂零：在 \(f_j\) 上由三次 Serre 关系保证，在 \(e_i\) 上先得到 \(h_i\)、\(f_i\) 的倍数，再变为零，在 \(e_j\)（\(j\ne i\)）上已为零，在 \(h_j,\mathrm d\) 上至多两步为零。导子法则使它在包络代数的每个有限乘积上也局部幂零。于是对任意 \(u\in U(\widehat{\mathfrak{sl}}_2)\)，可取 \(R\) 使 \((\operatorname{ad}f_i)^ru=0\) 对所有 \(r\ge R\) 成立。重排恒等式
	\[
		f_i^Nu v_\lambda
		=\sum_{r=0}^N\binom Nr
		(\operatorname{ad}f_i)^r(u)\,f_i^{N-r}v_\lambda
	\]
	表明，当 \(N\ge R+m_i+1\) 时，右侧每项都为零。因此 \(f_i\) 在全模上局部幂零。

	对于 \(e_i\)，只需使用权的范围。若 \(v\) 的权为 \(\lambda-r_0\alpha_0-r_1\alpha_1\)，则 \(e_i^Nv\) 的权将相应的 \(r_i\) 改为 \(r_i-N\)。当 \(N>r_i\) 时，这已不属于 \(\lambda-Q_+\)，所以 \(e_i^Nv=0\)。任意向量只有有限多个权分量，对这些分量取共同的上界即可。这证明 \(L(\lambda)\) 可积。
\end{proof}

支配整性在有限维半单情形保证简单最高权模有限维；在这里，它保证的是可积性。中心项直接说明这两个结论为何分开。
\begin{proposition}\label{b5:prop:G-positive-level-infinite}
	设 \(V\ne0\) 为复 Lie 代数 \(\widehat{\mathfrak{sl}}_2\) 的模，中心元 \(c\) 按标量 \(\ell\ne0\) 作用，则 \(V\) 无限维。可积简单最高权模 \(L(\lambda)\) 的级为零，当且仅当 \(\lambda=a\delta\)（\(a\in\mathbb C\)）；此时它是一维模，所有环代数元素及 \(c\) 作用为零，\(\mathrm d\) 按 \(a\) 作用。
\end{proposition}
\begin{proof}
	迹型满足 \(\operatorname{tr}(h^2)=2\)，所以
	\[
		[h\otimes t,h\otimes t^{-1}]=2c.
	\]
	若 \(V\) 有限维，在这个等式的表示上取迹，左边是两个有限维算子的交换子，迹为零；右边的迹为 \(2\ell\dim V\ne0\)，矛盾。

	在级为零的可积简单最高权模中，两个非负整数 \(m_0,m_1\) 之和为零，故它们都为零。上一条定理给出 \(f_0v_\lambda=f_1v_\lambda=0\)，而正生成元本来就零化 \(v_\lambda\)。由于这些元素生成全部环代数与中心，最高权生成元的张成空间已经是整个模。反过来，对任何 \(a\) 这样规定的一维作用都保持括号，简单且可积，其最高权为 \(a\delta\)。
\end{proof}

取级一的 \(L(\Lambda_0)\)，并记其最高权向量为 \(v\)。这时 \(m_0=1\)、\(m_1=0\)，所以
\[
	f_0^2v=0,\qquad f_1v=0,\qquad f_0v\ne0.
\]
最后一个不等式也可由 \(e_0f_0v=h_0v=v\) 直接看出。令 \(w=f_0v\)，则
\[
	e_1w=0,\qquad h_1w=2w.
\]
因此相对于零次 \(\mathfrak{sl}_2\)，它是一条最高权为 \(2\) 的根串的起点。局部幂零性与\eqref{b5:eq:G-simple-root-string}给出
\[
	w,\quad f_1w,\quad f_1^2w\ne0,\qquad f_1^3w=0.
\]
三个向量的 \(h_1\) 权分别为 \(2,0,-2\)，故线性无关；它们的 \(\mathrm d\) 特征值都为 \(-1\)。任何 \(\mathrm d\) 特征值为 \(-1\) 的下降字恰含一个 \(f_0\)，而 \(f_1v=0\) 使这个 \(f_0\) 右边不能再有 \(f_1\)。所以这三个向量张成整个 \(-1\) 次空间。\(\mathrm d\) 特征值为零的空间则只有 \(\mathbb Cv\)。若用非负整数 \(r\) 记录特征值 \(-r\)，头两层的维数因而为 \(1,3\)；将 PBW 基中的零次 \(f_1\) 排在最右边后，它在 \(v\) 上的正次幂全部消失；其余负根基元都有负的 \(\mathrm d\) 次数，每个固定总次数只允许有限多个 PBW 单项式。因此后续层也各自有限维，却不会在某一层后全都消失。

% !TeX root = ../../Book5.tex
% 纲要标题：### G.4 Weyl–Kac 公式与小秩例子
% 纲要位置：工作记录/计划纲要.md（2026-10-20 新增骨架）。
% 写作状态：2026-10-20 新增骨架；正文、证明、算例与习题均待撰写。
% 写作范围：以可对称化及仿射型为范围，展开环代数和仿射 sl2 小模型，Weyl–Kac 公式标明证明范围；不展开 VOA 与 Lie 超代数。
% #### G.4.1 形式特征标与完成
% #### G.4.2 Weyl–Kac 特征标公式
% #### G.4.3 仿射秩一与秩二的例子

\section{Weyl–Kac 公式与小秩例子}
\label{b5:sec:G04}

待撰写。

\subsection{形式特征标与完成}
\label{b5:subsec:G04:01}

待撰写。

\subsection{Weyl–Kac 特征标公式}
\label{b5:subsec:G04:02}

待撰写。

\subsection{仿射秩一与秩二的例子}
\label{b5:subsec:G04:03}

待撰写。

% 新增结构与纲要同步；以下条目均为待撰写计划。
% * 状态：待撰写；当前仅有结构占位。
% ---


\begin{chaptersummary}
	广义 Cartan 矩阵的奇异性同时留下中心方向和额外的 Cartan 方向；在仿射环模型中，前者由中心元 \(c\) 实现，后者由次数导子 \(\mathrm d\) 实现。实根是有限根加上整数倍 \(\delta\)，虚根则是非零整数倍 \(\delta\)，重数等于原有限根系的秩。两类根在形式特征标的分母中分别出现。
	\ProseParagraphBreak
	对 \(\widehat{\mathfrak{sl}}_2\)，最高权在两个简单余根上的取值必须同时为非负整数，简单最高权模才可积；两者之和是级。非零级迫使表示无限维，却不妨碍每个简单根方向局部幂零，也不妨碍每个固定次数层有限维。Weyl–Kac 公式用反射轨道的交错和扣除 Verma 模中的关系，级一模 \(L(\Lambda_0)\) 的前三层维数为 \(1,3,4\)，其中前两层亦由具体向量计算得到。形式完成保证逐权计算有意义，合并变量或求维数时还须检查新的系数是否有限。
\end{chaptersummary}

\BookExercises


\begin{exercise}\label{b5:ex:G-symmetrizing-cycle}
	设 \(A=(a_{ij})\) 为不可分解广义 Cartan 矩阵。证明：\(A\) 可对称化，当且仅当其连接图中每条闭路 \(i_0,i_1,\ldots,i_m=i_0\) 都满足
	\[
		\prod_{k=0}^{m-1}\frac{a_{i_k i_{k+1}}}{a_{i_{k+1}i_k}}=1.
	\]
	这里每相邻两个顶点相连，闭路允许重复顶点。用这一判据说明
	\[
		A=\begin{pmatrix}2&-1&-1\\-2&2&-1\\-1&-2&2\end{pmatrix}
	\]
	不是可对称化矩阵，并说明连接图为树时为何总可对称化。
\end{exercise}
\begin{solution}
	若 \(D=\operatorname{diag}(d_i)\) 为正对角对称化矩阵，则每条边上有 \(a_{ij}/a_{ji}=d_j/d_i\)。沿闭路相乘，各 \(d_i\) 在分子、分母中抵消，乘积为一。

	反过来，固定顶点 \(i_*\)，令 \(d_{i_*}=1\)。对从 \(i_*\) 到 \(j\) 的路径，定义 \(d_j\) 为沿路径的比值 \(a_{uv}/a_{vu}\) 的乘积。每个因子为正。两条路径相接成一条闭路；反向走边时比值取倒数，因此闭路条件说明两条路径的乘积相同，\(d_j\) 良定义。对相连的 \(i,j\)，用通向 \(i\) 的路径再添一条边，便有 \(d_j=d_i a_{ij}/a_{ji}\)，即 \(d_i a_{ij}=d_j a_{ji}\)。不相连处两个矩阵元都为零，同样满足等式，所以 \(DA\) 对称。所有 \(d_i\) 都是正有理数，清分母可得到正整数对称化矩阵。

	题中三角闭路 \(1,2,3,1\) 的比值乘积为 \((1/2)(1/2)\cdot1=1/4\ne1\)，故不能对称化。若图为树，则两点间有唯一不重复顶点的路径，可以直接用这条路径定义 \(d_j\)。每条闭路都能逐次删除一条往返边，两方向的比值互为倒数，因此闭路条件自动成立。可对称化的障碍由不同路径的相容性产生，单条边的正比值本身并不足以构成障碍。
\end{solution}

\begin{exercise}\label{b5:ex:G-three-vertex-chain}
	设 \(a,b,c,d\) 为正整数，考虑
	\[
		A=\begin{pmatrix}2&-a&0\\-b&2&-c\\0&-d&2\end{pmatrix}.
	\]
	求一个正整数对称化矩阵，并证明：\(A\) 为有限型当且仅当 \(ab+cd<4\)，为仿射型当且仅当 \(ab+cd=4\)。在仿射情形，分别给出 \(A\) 与 \(A^{\mathsf T}\) 的一个全正零向量，并据此写出实现中的一个非零中心元。
\end{exercise}
\begin{solution}
	取 \(D=\operatorname{diag}(bd,ad,ac)\)，则
	\[
		DA=\begin{pmatrix}2bd&-abd&0\\-abd&2ad&-acd\\0&-acd&2ac\end{pmatrix}.
	\]
	它的二次型可以配方为
	\[
		\begin{aligned}
			B(x,y,z)
			&=2bd\left(x-\frac a2y\right)^2
			+2ac\left(z-\frac d2y\right)^2
			+\frac{ad(4-ab-cd)}2y^2.
		\end{aligned}
	\]
	三个新坐标 \(x-ay/2,z-dy/2,y\) 由原坐标经可逆线性变换得到，前两个系数严格为正。因此最后一个系数为正时正定，为零时半正定且余秩一，为负时同时取正、负值。连接图是一条三顶点链，故不可分解，所需判别随之成立。

	若 \(ab+cd=4\)，直接相乘得
	\[
		A\begin{pmatrix}a\\2\\d\end{pmatrix}=0,
		\qquad
		A^{\mathsf T}\begin{pmatrix}b\\2\\c\end{pmatrix}=0.
	\]
	所以在任意复实现中，\(a\alpha_1+2\alpha_2+d\alpha_3\) 零化全部 \(h_i\)，而 \(bh_1+2h_2+ch_3\) 被全部 \(\alpha_i\) 零化。后一元素由\cref{b5:prop:G-central-directions}成为非零中心元；它非零是因为 \(h_i\) 独立。前一泛函也非零，因为简单根独立，只是在 \(\mathfrak h'\) 上为零。左右核通常给出不同的系数，不能在非对称矩阵中直接交换两组数值。
\end{solution}

\begin{exercise}\label{b5:ex:G-two-affine-components}
	令 \(A_*=\begin{pmatrix}2&-2\\-2&2\end{pmatrix}\)，并令 \(A=\operatorname{diag}(A_*,A_*)\)。为两块分别取\cref{b5:def:G-realization}之后的三维实现，其简单余根、额外方向记为 \(h_0^{(j)},h_1^{(j)},\mathrm d_j\)，其中 \(j=1,2\)。证明两份实现的直和是 \(A\) 的实现，并证明相应的 Kac--Moody 代数为
	\[
		\mathfrak g(A)\cong\mathfrak g(A_*)\oplus\mathfrak g(A_*).
	\]
	求 \(\mathfrak g(A)\) 的中心和交换化，并解释为什么 \(A\) 整体不是本节定义下的仿射型。
\end{exercise}
\begin{solution}
	矩阵 \(A\) 的阶数为 \(4\)、秩为 \(2\)，所以实现所需的维数为 \(2\cdot4-2=6\)。在直和中，把每一块的简单根延拓为在另一块上为零的泛函。四个简单根与四个简单余根各自独立，配对矩阵恰为 \(A\)，因此这确是一个实现。

	来自不同块的正生成元满足 \([e_i^{(1)},e_k^{(2)}]=0\)，因为交叉矩阵元为零，Serre 指数为一；负生成元同样交换。正负交叉括号、两个实现空间之间的括号以及一个实现空间对另一块升降生成元的作用也都为零。因此两块生成的子代数彼此交换。两份展示给出一个从 \(\mathfrak g(A_*)\oplus\mathfrak g(A_*)\) 到 \(\mathfrak g(A)\) 的满同态；反向把每个生成元送入相应直和因子，全部关系仍成立，故也给出同态。两个复合都在全部生成元上为恒等，于是得到所述同构。

	记 \(c_j=h_0^{(j)}+h_1^{(j)}\)。每个因子的中心为 \(\mathbb Cc_j\)，所以整体中心为 \(\mathbb Cc_1\oplus\mathbb Cc_2\)。交换化则以 \(\mathrm d_1,\mathrm d_2\) 的像为基，是二维交换 Lie 代数。矩阵 \(A\) 虽半正定，却可分解且核为二维；它是两个仿射型块的直和，而不是一个不可分解仿射型矩阵。两个中心方向和两个额外方向也分别保留，不能用一个共同的 \(c\) 或 \(\mathrm d\) 替代。
\end{solution}

\begin{exercise}\label{b5:ex:G-positive-affine-kernel}
	设 \(A\) 为 \(r\) 阶仿射型广义 Cartan 矩阵，\(D\) 为正对角对称化矩阵，\(B=DA\)。证明 \(\ker A\) 由一个所有坐标均严格为正的向量张成，而且这个向量可以取成正整数向量。再证明 \(B\) 的每个非空真主子矩阵都正定。这里主子矩阵使用同一指标集选取行与列。
\end{exercise}
\begin{solution}
	取非零实向量 \(v\in\ker B\)，令 \(w_i=|v_i|\)。由于 \(b_{ij}\le0\) 对 \(i\ne j\) 成立，有
	\[
		w^{\mathsf T}Bw
		=\sum_i b_{ii}v_i^2+2\sum_{i<j}b_{ij}|v_iv_j|
		\le v^{\mathsf T}Bv=0.
	\]
	半正定性迫使左端等于零。由实对称矩阵的正交对角化，半正定型取零的向量属于矩阵的核，故 \(Bw=0\)。若某个 \(w_i=0\)，则该行方程化为 \(\sum_{j\ne i}b_{ij}w_j=0\)。每个求和项都非正，所以只要 \(b_{ij}<0\)，就有 \(w_j=0\)。连接图连通，这个零值会沿全部边传播，使 \(w=0\)，与 \(v\ne0\) 矛盾。因此每个 \(w_i>0\)。

	由于 \(\ker B=\ker A\) 为一维，它恰为 \(\mathbb Rw\)。整数矩阵 \(A\) 的秩为 \(r-1\)，有理数上的行化简给出一个非零有理零向量；它与 \(w\) 成比例，故可以整体变号成为全正向量，再清分母成为正整数向量。

	设 \(I\) 是非空真指标集，\(B_I\) 是相应主子矩阵。对 \(x\in\mathbb R^I\)，在其余坐标补零得到 \(\widetilde x\in\mathbb R^r\)，则 \(x^{\mathsf T}B_Ix=\widetilde x^{\mathsf T}B\widetilde x\ge0\)。若它为零，则 \(\widetilde x\in\ker B=\mathbb Rw\)。但 \(\widetilde x\) 至少有一个坐标为零，而 \(w\) 的每个坐标都正，所以 \(\widetilde x=0\)，从而 \(x=0\)。这证明 \(B_I\) 正定。
\end{solution}



\begin{exercise}\label{b5:ex:G-evaluation-no-degree}
	设 \(\mathfrak g\) 为有限维复单 Lie 代数，\(a\in\mathbb C^\times\)。
	\begin{enumerate}
		\item 证明求值映射 \(\operatorname{ev}_a:L\mathfrak g\rightarrow\mathfrak g\)、\(x\otimes f\mapsto f(a)x\) 为满射 Lie 同态，核为 \(\mathfrak g\otimes(t-a)\mathbb C[t,t^{-1}]\)。
		\item 设 \(\rho:\mathfrak g\rightarrow\mathfrak{gl}(V)\) 为非零表示。令 \(x\otimes t^m\) 在 \(V\) 上作用为 \(a^m\rho(x)\)，并令中心元 \(c\) 作用为零。证明这给出 \(\widetilde{\mathfrak g}\) 的表示，但不存在 \(\mathrm d\) 在同一空间上的作用，使它延拓成 \(\widehat{\mathfrak g}\) 的表示。
	\end{enumerate}
\end{exercise}
\begin{solution}
	1. 求值保持乘法，因而
	\[
		\operatorname{ev}_a([x\otimes f,y\otimes g])
		=f(a)g(a)[x,y]
		=[\operatorname{ev}_a(x\otimes f),\operatorname{ev}_a(y\otimes g)].
	\]
	常数环元给出满射。取 \(\mathfrak g\) 的基 \(x_1,\ldots,x_r\)，则 \(\sum_i x_i\otimes f_i\) 落入核当且仅当各 \(f_i(a)=0\)。对每个 Laurent 多项式，乘以某个 \(t^N\) 后成为普通多项式；因 \(a\ne0\)，求值为零当且仅当这个多项式被 \(t-a\) 整除，而 \(t^N\) 在 Laurent 多项式环中可逆。故每个 \(f_i\) 都属于 \((t-a)\)，得到所述核。

	2. 中心作用为零以后，中心修正项也作用为零，所以第一问的同态与 \(\rho\) 复合就给出表示。若某个算子 \(D\) 能代表 \(\mathrm d\)，则由 \([\mathrm d,x\otimes1]=0\) 得
	\[
		[D,\rho(x)]=0\qquad(x\in\mathfrak g).
	\]
	但由 \([\mathrm d,x\otimes t]=x\otimes t\) 又应有
	\[
		[D,a\rho(x)]=a\rho(x).
	\]
	左边为零，且 \(a\ne0\)，于是 \(\rho=0\)，矛盾。求值把次数压成数值 \(a^m\)，因此一般不再容许一个同时测量这些次数的导子。
\end{solution}

\begin{exercise}\label{b5:ex:G-loop-coordinate-changes}
	设 \(\widehat{\mathfrak g}\) 为有限维复单 Lie 代数的未扭仿射代数，括号沿用\eqref{b5:eq:G-affine-bracket}。
	\begin{enumerate}
		\item 对 \(a\in\mathbb C^\times\)，证明 \(u_m\mapsto a^m u_m\)、\(c\mapsto c\)、\(\mathrm d\mapsto\mathrm d\) 是自同构。
		\item 证明 \(u_m\mapsto u_{-m}\)、\(c\mapsto-c\)、\(\mathrm d\mapsto-\mathrm d\) 是一个二阶自同构。若保持 \(c\) 不变，为什么 \(u_m\mapsto u_{-m}\) 不能给出中心扩张的自同构？
	\end{enumerate}
\end{exercise}
\begin{solution}
	1. 两个变换后齐次元素的括号为
	\[
		[a^m u_m,a^n v_n]
		=a^{m+n}[u,v]_{m+n}+a^{m+n}m\delta_{m,-n}(u,v)c.
	\]
	只要中心项非零，就有 \(m+n=0\) 和 \(a^{m+n}=1\)，因此这等于原括号的像。次数导子的关系也保持，因为 \([\mathrm d,a^m u_m]=m a^m u_m\)。把 \(a\) 换成 \(a^{-1}\) 得到逆映射。

	2. 记所给映射为 \(\iota\)。一方面，
	\[
		[\iota(u_m),\iota(v_n)]
		=[u,v]_{-m-n}-m\delta_{m,-n}(u,v)c;
	\]
	另一方面，这正是 \(\iota([u_m,v_n])\)，因为 \(\iota(c)=-c\)。此外
	\[
		[\iota(\mathrm d),\iota(u_m)]
		=[-\mathrm d,u_{-m}]=m u_{-m}=\iota(mu_m).
	\]
	其余中心关系自动保持，且两次应用回到原元素，故 \(\iota^2=1\)。如果保持 \(c\) 不变，取 \((h,h)\ne0\)，则 \([h_1,h_{-1}]=(h,h)c\) 的像应为 \((h,h)c\)，但两个像的括号为 \([h_{-1},h_1]=-(h,h)c\)，矛盾。变量反演不仅反转次数，也反转这份中心修正项的符号。
\end{solution}

\begin{exercise}\label{b5:ex:G-heisenberg-polynomials}
	在 \(\widehat{\mathfrak{sl}}_2\) 中，令 \(a_n=h\otimes t^n\)，并令 \(\mathfrak s\) 为 \(c,\mathrm d\) 及全部 \(a_n\)（\(n\ne0\)）的线性生成空间。对固定 \(k\in\mathbb C\)，在只含有限多个变量的普通多项式所成空间 \(P=\mathbb C[z_1,z_2,\ldots]\) 上规定
	\[
		\begin{aligned}
			a_{-n}\cdot F&=z_nF,& a_n\cdot F&=2nk\frac{\partial F}{\partial z_n},\\
			c\cdot F&=kF,& \mathrm d\cdot F&=-\sum_{n\ge1}n z_n\frac{\partial F}{\partial z_n}
		\end{aligned}
		\qquad(n\ge1).
	\]
	证明 \(\mathfrak s\) 是 Lie 子代数，所列公式给出它在 \(P\) 上的表示；并证明当 \(k\ne0\) 时，此表示限制到 \(c,a_n\) 生成的子代数上仍不可约。
\end{exercise}
\begin{solution}
	由\eqref{b5:eq:G-affine-sl2-brackets}，有
	\[
		[a_m,a_n]=2m\delta_{m,-n}c,\qquad
		[\mathrm d,a_n]=na_n,
	\]
	所以题中空间在括号下封闭。每个 \(F\) 只涉及有限多个 \(z_n\)，故 \(\mathrm d\) 的求和在每个多项式上实际都是有限和。乘法算子彼此交换，偏导算子彼此交换，而
	\[
		\left[2mk\frac{\partial}{\partial z_m},z_n\right]
		=2mk\delta_{mn}I,
	\]
	正好实现相反次数的中心括号。带权 Euler 算子 \(E=\sum_{n\ge1}n z_n\partial/\partial z_n\) 满足 \([E,z_n]=nz_n\)、\([E,\partial/\partial z_n]=-n\partial/\partial z_n\)；取 \(\mathrm d=-E\)，就得到两个次数关系。

	设 \(k\ne0\)，\(W\subseteq P\) 为非零子模，取非零 \(F\in W\)。选取 \(F\) 中一个全次数最大的单项式 \(z_1^{b_1}\cdots z_r^{b_r}\)，用对应的偏导 \(\partial_1^{b_1}\cdots\partial_r^{b_r}\) 作用。低于该全次数的单项式都被零化；同一全次数的其他单项式若不同于所选项，就至少有一个指数小于相应 \(b_i\)，也被零化。因此结果是一个非零常数。由于 \(2nk\ne0\)，每个所用偏导都是 \(a_n\) 作用的非零标量倍数，所以 \(1\in W\)。再用 \(a_{-n}\) 逐次乘上变量，就得到全部单项式，故 \(W=P\)。
\end{solution}

\begin{exercise}\label{b5:ex:G-finite-affine-representations}
	设 \(\mathfrak g\) 为有限维复单 Lie 代数，\(\widehat{\mathfrak g}\) 为其未扭仿射 Lie 代数。证明每个有限维表示 \(\rho:\widehat{\mathfrak g}\rightarrow\mathfrak{gl}(V)\) 都零化 \(\widetilde{\mathfrak g}=L\mathfrak g\oplus\mathbb Cc\)。因此这类表示只剩 \(\mathrm d\) 的任意一个线性算子作用。
\end{exercise}
\begin{solution}
	令 \(D=\rho(\mathrm d)\)。有限维空间 \(\operatorname{End}_{\mathbb C}(V)\) 上的线性算子 \(A\mapsto[D,A]\) 只有有限多个特征值。关系
	\[
		[D,\rho(x_m)]=m\rho(x_m)
	\]
	说明，只要整数 \(m\) 不在这个有限特征值集合内，就有 \(\rho(x_m)=0\) 对所有 \(x\in\mathfrak g\) 成立。取一个正整数 \(m\) 充分大，使 \(m,-m,m+1,-m-1\) 都不在该集合内。对任意 \(x,y\in\mathfrak g\)，相反次数括号分别给出
	\[
		0=\rho([x,y]_0)+m(x,y)\rho(c),\qquad
		0=\rho([x,y]_0)+(m+1)(x,y)\rho(c).
	\]
	相减并选取 \((x,y)\ne0\)，得到 \(\rho(c)=0\)。再由 \([\mathfrak g,\mathfrak g]=\mathfrak g\) 得 \(\rho(\mathfrak g\otimes1)=0\)。最后，对任意 \(n\in\mathbb Z\)，关系 \([x_0,y_n]=[x,y]_n\) 说明 \(\rho([x,y]_n)=0\)；仍用有限维单代数的完美性，便得 \(\rho(\mathfrak g\otimes t^n)=0\)。这证明整个 \(\widetilde{\mathfrak g}\) 作用为零。反过来，令该理想作用为零，而让 \(\mathrm d\) 作用为任意一个算子，显然满足所有括号关系。
\end{solution}



\begin{exercise}\label{b5:ex:G-affine-node-symmetry}
	在以迹型构造的复仿射 Lie 代数 \(\widehat{\mathfrak{sl}}_2\) 中，对每个 \(r\in\mathbb Z\) 定义线性映射
	\[
		\begin{aligned}
			\sigma(e\otimes t^r)&=f\otimes t^{r+1},&
			\sigma(f\otimes t^r)&=e\otimes t^{r-1},\\
			\sigma(h\otimes t^r)&=-h\otimes t^r+\delta_{r0}c,&
			\sigma(c)&=c,\\
			\sigma(\mathrm d)&=\mathrm d+\frac12h-\frac14c.
		\end{aligned}
	\]
	证明它是交换两组 \(e_i,f_i,h_i\) 的 Lie 代数对合。对任意最高权模 \(V\)，令 \(V^\sigma\) 具有相同向量空间与新作用 \(x\cdot_\sigma v=\sigma(x)v\)。求 \(L(\lambda)^\sigma\) 的最高权，并证明 \(L(\Lambda_0)^\sigma\cong L(\Lambda_1-\dfrac14\delta)\)。
\end{exercise}
\begin{solution}
	最需要核对的是中心项。记 \(E_r=e\otimes t^r\)、\(F_r=f\otimes t^r\)、\(H_r=h\otimes t^r\)，则
	\[
		\begin{aligned}
			[\sigma(E_r),\sigma(F_s)]
			&=[F_{r+1},E_{s-1}]
			=-H_{r+s}+(r+1)\delta_{r,-s}c,\\
			\sigma([E_r,F_s])
			&=\sigma(H_{r+s}+r\delta_{r,-s}c)
			=-H_{r+s}+(r+1)\delta_{r,-s}c.
		\end{aligned}
	\]
	两式相同。\([H_r,H_s]=2r\delta_{r,-s}c\) 也被保持，因为 \(\sigma(H_r)\) 与 \(-H_r\) 仅差一个中心项。其余环模态的括号由
	\[
		[-H_r,F_{s+1}]=2F_{r+s+1},\qquad
		[-H_r,E_{s-1}]=-2E_{r+s-1}
	\]
	及同种 \(e\)、同种 \(f\) 之间的零括号得到。最后
	\[
		[\mathrm d+\frac12h-\frac14c,F_{r+1}]=r F_{r+1},\qquad
		[\mathrm d+\frac12h-\frac14c,E_{r-1}]=r E_{r-1},
	\]
	而与 \(\sigma(H_r)\) 的括号为 \(-r H_r=\sigma(rH_r)\)。因此 \(\sigma\) 保持全部括号。

	在环模态及 \(c\) 上连续作用两次可直接得到恒等；在 \(\mathrm d\) 上有
	\[
		\sigma^2(\mathrm d)
		=\bigl(\mathrm d+\frac12h-\frac14c\bigr)
		+\frac12(-h+c)-\frac14c=\mathrm d.
	\]
	所以它是对合。代入简单生成元便有 \(\sigma(e_0)=e_1\)、\(\sigma(e_1)=e_0\)，\(f_i,h_i\) 同样交换。

	\(\sigma\) 保持 \(\mathfrak n_+\) 和 \(\widehat{\mathfrak h}\)，所以原最高权向量在扭曲模中仍为最高权向量，其权为 \(\lambda\circ\sigma\)。若 \(\lambda=m_0\Lambda_0+m_1\Lambda_1+a\delta\)，则
	\[
		\lambda\circ\sigma
		=m_1\Lambda_0+m_0\Lambda_1
		+\left(a+\frac{m_1-m_0}{4}\right)\delta.
	\]
	自同构扭曲保持子模、简单性及生成性，因此 \(L(\lambda)^\sigma\cong L(\lambda\circ\sigma)\)。取 \(\lambda=\Lambda_0\) 即得所需同构。中心项与次数导子中的修正同时影响这个结论，不能只交换两个基本权而保持 \(\mathrm d\) 不变。
\end{solution}

\begin{exercise}\label{b5:ex:G-degree-shift}
	设 \(\mathbb C_a\) 是复仿射 Lie 代数 \(\widehat{\mathfrak{sl}}_2\) 的一维模，全部环模态和 \(c\) 作用为零，\(\mathrm d\) 按 \(a\in\mathbb C\) 作用。证明
	\[
		L(\lambda)\otimes\mathbb C_a\cong L(\lambda+a\delta).
	\]
	说明这些模的级不变，限制到导出代数后同构；但对不同的 \(a\)，它们作为整个仿射代数的最高权模互不同构。
\end{exercise}
\begin{solution}
	在 \(L(\lambda)\otimes\mathbb C_a\) 上，环模态和 \(c\) 的作用与 \(L(\lambda)\) 相同，\(\mathrm d\) 的作用增加了 \(aI\)。因此两个模有相同的子模；原模简单，张量后的模也简单。最高权向量 \(v_\lambda\otimes1\) 的权在 \(h,c\) 上不变，在 \(\mathrm d\) 上增加 \(a\)，所以其权为 \(\lambda+a\delta\)。它仍生成全模，由简单最高权模的唯一性得到所示同构。

	因为 \(\delta(c)=0\)，级仍为 \(\lambda(c)\)。导出代数由环模态和 \(c\) 组成，不含 \(\mathrm d\)，所以限制后的作用完全相同。对 \(a\ne b\)，最高权 \(\lambda+a\delta\ne\lambda+b\delta\)；\cref{b5:prop:G-affine-verma}给出整个仿射代数模之间不存在同构。
\end{solution}

\begin{exercise}\label{b5:ex:G-integrable-tensor}
	证明两个复仿射 \(\widehat{\mathfrak{sl}}_2\) 可积模的张量积仍可积；若它们的级分别为 \(\ell_1,\ell_2\)，则张量积的级为 \(\ell_1+\ell_2\)。再取 \(V=L(\Lambda_0)\)，以 \(v\) 表示最高权生成元，证明
	\[
		w=f_0v\otimes v-v\otimes f_0v
	\]
	是 \(V\otimes V\) 中非零的最高权向量。求它的权，并据此证明 \(V\otimes V\) 不简单。
\end{exercise}
\begin{solution}
	张量积中的纯张量权为两个因子权之和，故仍按权直和分解。固定 \(i\)，\(e_i\) 在张量积上的作用是 \(e_i\otimes1+1\otimes e_i\)。对 \(u\otimes z\)，若 \(e_i^ru=0\)、\(e_i^sz=0\)，二项式展开表明 \(e_i^{r+s}(u\otimes z)=0\)。\(f_i\) 同理，任意张量又是有限多个纯张量之和，故二者局部幂零。中心元的作用为 \(c\otimes1+1\otimes c\)，即标量 \(\ell_1+\ell_2\)。

	由正文算例，\(f_0v\ne0\)、\(e_0f_0v=v\)，而 \(v\) 与 \(f_0v\) 权不同，线性无关。因此两个纯张量 \(f_0v\otimes v\)、\(v\otimes f_0v\) 线性无关，\(w\ne0\)。又有
	\[
		e_0w=v\otimes v-v\otimes v=0,
	\]
	且 \(e_1v=e_1f_0v=0\)，故 \(e_1w=0\)。它的权为 \(2\Lambda_0-\alpha_0\)，所以 \(w\) 生成一个最高权子模，其所有权都在 \(2\Lambda_0-\alpha_0-Q_+\) 中。这个子模不能包含非零向量 \(v\otimes v\)，后者权为 \(2\Lambda_0\)。于是它是非零真子模，证明张量积不简单。
\end{solution}

\begin{exercise}\label{b5:ex:G-level-one-second-module}
	设 \(V=L(\Lambda_1)\)，\(v\) 为其最高权生成元。证明 \(\mathrm d\) 特征值为 \(0\) 的空间有基 \(v,f_1v\)，特征值为 \(-1\) 的空间有基
	\[
		f_0f_1v,\qquad f_1f_0f_1v.
	\]
	计算这四个向量的 \(h_1\) 权，并说明它们与 \(L(\Lambda_0)\) 中头两层的维数并不相同。
\end{exercise}
\begin{solution}
	此时 \(\Lambda_1(h_0)=0\)、\(\Lambda_1(h_1)=1\)，可积最高权判据给出 \(f_0v=0\)、\(f_1^2v=0\)，并且 \(f_1v\ne0\)。所有 \(\mathrm d\) 特征值为零的下降字都只含 \(f_1\)，所以该空间恰为 \(\mathbb Cv\oplus\mathbb Cf_1v\)，两个 \(h_1\) 权为 \(1,-1\)。

	令 \(u=f_0f_1v\)。利用 \([e_0,f_1]=0\) 及 \(e_0v=0\)，得到
	\[
		e_0u=h_0f_1v=2f_1v\ne0,
	\]
	故 \(u\ne0\)。此外
	\[
		e_1u=f_0e_1f_1v=f_0v=0,
		\qquad h_1u=u.
	\]
	它是零次 \(\mathfrak{sl}_2\) 中最高权为 \(1\) 的向量。局部幂零性与根串公式说明 \(f_1u\ne0\)、\(f_1^2u=0\)，两向量的 \(h_1\) 权为 \(1,-1\)。

	特征值 \(-1\) 的下降字恰含一个 \(f_0\)，故形如 \(f_1^rf_0f_1^sv\)。当 \(s=0\) 时，由 \(f_0v=0\) 得零；当 \(s\ge2\) 时，由 \(f_1^2v=0\) 得零；剩下 \(s=1\)，便是 \(f_1^ru\)。因此 \(u,f_1u\) 张成该层，且由互异的 \(h_1\) 权线性无关。头两层维数为 \(2,2\)，与 \(L(\Lambda_0)\) 的 \(1,3\) 不同。两者同为级一，但级只记录中心作用，并不决定最高权模。
\end{solution}



\begin{exercise}\label{b5:ex:G-completion-boundary}
	设 \(\alpha_0,\alpha_1\) 线性无关。在负根方向的完成中计算
	\[
		\frac1{(1-e^{-\alpha_0})(1-e^{-\alpha_1})(1-e^{-(\alpha_0+\alpha_1)})}
	\]
	里 \(e^{-a\alpha_0-b\alpha_1}\) 的系数，其中 \(a,b\) 为非负整数。说明为什么 \(\sum_{n\ge0}e^{n\alpha_1}\) 与 \(\sum_{n\ge0}e^{-n\alpha_1}\) 不能作为两个任意形式和直接相乘。再说明即使负根方向的乘积良定义，也不能总把 \(e^{-\alpha_1}\) 代成一。
\end{exercise}
\begin{solution}
	三类因子的次数分别记为 \(u,v,w\)。所求系数等于非负整数解 \(u+w=a\)、\(v+w=b\) 的个数。\(w\) 可取 \(0,\ldots,\min(a,b)\)，所以系数为 \(\min(a,b)+1\)。这也给出逐权计算只有有限项的具体例子。

	第二个乘积的常数项会收到所有 \(e^{n\alpha_1}e^{-n\alpha_1}\) 的贡献，成为无限个一之和，因而没有整数系数意义。第一个因子的支集也不在有限多个负根锥的平移中。最后，合法的几何级数 \(\sum_{n\ge0}e^{-n\alpha_1}\) 在形式变量代成一后就遇到同一问题；形式完成只保证原有权系数的有限性，不保证合并无限多个权仍可求和。
\end{solution}

\begin{exercise}\label{b5:ex:G-imaginary-partitions}
	设 \(r\) 为正整数，考虑只记录虚根 PBW 因子的形式级数
	\[
		P_r(q)=\prod_{n\ge1}(1-q^n)^{-r}.
	\]
	把其系数解释为带 \(r\) 种颜色的分拆数，求 \(q,q^2,q^3\) 的系数。比较 \(r=1\) 与 \(r=2\)，说明把 \(A_2^{(1)}\) 的虚根重数误写成一会怎样改变计数。这一乘积是否就是简单可积最高权模的完整特征标？
\end{exercise}
\begin{solution}
	对每个正整数 \(n\)，有 \(r\) 个可重复使用的次数为 \(n\) 的基向量，分别看成 \(r\) 种颜色的大小 \(n\) 的部分。因此 \(q\) 的系数为 \(r\)。次数二可由一个大小二的部分或两个大小一的部分得到，系数为 \(r+\binom{r+1}{2}\)。次数三的大小类型为 \(3\)、\(2+1\)、\(1+1+1\)，系数为
	\[
		r+r^2+\binom{r+2}{3}.
	\]
	\(r=1\) 时前三个系数为 \(1,2,3\)，\(r=2\) 时为 \(2,5,10\)。虚根的重数因此已经影响最低几次计数。

	这个乘积只记录负虚根基向量的 PBW 因子。完整 Verma 特征标还包含所有负实根因子和最高权平移；简单最高权模又是 Verma 模的商，其关系要由 Weyl–Kac 分子等信息扣除。因此不能把 \(P_r(q)\) 本身当作该简单模的完整特征标。
\end{solution}

\begin{exercise}\label{b5:ex:G-level-vacuum-grades}
	对 \(\widehat{\mathfrak{sl}}_2\) 及整数 \(k\ge1\)，以 \(x=e^{\alpha_1}\)、\(q=e^{-\delta}\) 写出 \(e^{-k\Lambda_0}\operatorname{ch}L(k\Lambda_0)\) 的 Weyl–Kac 分子。证明分子除去 \(1-x^{-1}\) 后，首个非常数 \(q\) 次项在 \(q^{k+1}\) 出现。由此求 \(k\ge2\) 时 \(\mathrm d\) 特征值 \(0,-1,-2\) 三层的维数，并与级一情形比较。
\end{exercise}
\begin{solution}
	\(k\Lambda_0+\rho\) 的级为 \(K=k+2\)，且 \((z,a)=(1,0)\)。\eqref{b5:eq:G-affine-translation}给出归一化分子
	\[
		N_k=\sum_{j\in\mathbb Z}\bigl(q^{Kj^2+j}x^{Kj}
		-q^{Kj^2-j}x^{Kj-1}\bigr).
	\]
	常数项为 \(1-x^{-1}\)。\(j=1\) 的负项与 \(j=-1\) 的正项在次数 \(K-1=k+1\) 首次出现，合为 \((x^{-K}-x^{K-1})q^{K-1}\)，其系数非零。\(|j|\ge2\) 时次数至少为 \(4K-2>K-1\)，所以没有更低的项。除去 \(1-x^{-1}\) 不改变 \(q\) 次数，首项仍然非零。

	若 \(k\ge2\)，该修正从次数三或更高才出现。因此前两次系数只是\eqref{b5:eq:G-sl2-denominator}中无限乘积的逆所给的
	\[
		x+1+x^{-1},\qquad x^2+2x+3+2x^{-1}+x^{-2}.
	\]
	三层维数为 \(1,3,9\)。级一时分子已在次数二扣除 \(x^2+x+1+x^{-1}+x^{-2}\)，三层维数改为 \(1,3,4\)。这说明同一组根和同一分母并不决定简单模的所有权重数；级会改变关系首次出现的次数。
\end{solution}

\begin{exercise}\label{b5:ex:G-horizontal-character}
	设 \(\lambda=m\Lambda_0+n\Lambda_1\)，其中 \(m,n\in\mathbb Z_{\ge0}\)，且 \(\lambda(\mathrm d)=0\)。记 \(x=e^{\alpha_1}\)、\(q=e^{-\delta}\)。利用仿射反射公式，求 \(e^{-\lambda}\operatorname{ch}L(\lambda)\) 的 \(q^0\) 项，并由此辨认零次层作为 \(\mathfrak{sl}_2\) 模的特征标。另设 \(b\in\mathbb C\)，证明若把 Weyl–Kac 公式中的 \(\rho\) 改成 \(\rho+b\delta\)，公式不变。
\end{exercise}
\begin{solution}
	令 \(K=m+n+2\)、\(z=n+1\)。\(\lambda+\rho\) 的坐标为 \((z,0)\)。交错分子除去 \(e^\lambda\) 后为
	\[
		\sum_{j\in\mathbb Z}\left(q^{Kj^2+jz}x^{Kj}
		-q^{Kj^2-jz}x^{Kj-z}\right).
	\]
	由于 \(0<z<K\)，当 \(j\ne0\) 时两种 \(q\) 指数都严格为正。故常数项只有 \(1-x^{-(n+1)}\)；分母常数项为 \(1-x^{-1}\)。其商为
	\[
		1+x^{-1}+\cdots+x^{-n}.
	\]
	乘回最高权后，零次层的 \(h\) 权为 \(n,n-2,\ldots,-n\)，每个重数为一。这是有限维简单 \(\mathfrak{sl}_2\) 模 \(L(n)\) 的特征标，维数为 \(n+1\)，与最高权的有限节点系数一致。

	因为 \(\delta(h_i)=0\)，每个反射 \(s_i\) 都固定 \(\delta\)。因此
	\[
		w(\lambda+\rho+b\delta)-(\rho+b\delta)
		=w(\lambda+\rho)-\rho.
	\]
	分子逐项不变，分母也与 \(\rho\) 无关。选 \(\rho(\mathrm d)=0\) 只是固定代表，并不改变公式。
\end{solution}
